\section{선형시불변 시스템}

전기적 시스템과 기계적 시스템을 같은 언어로 비교하려면 먼저 선형시불변 시스템(Linear Time-Invariant system, LTI system)을 이해해야 한다.
여기서 시스템(system)은 입력 신호를 받아 출력 신호를 만들어 내는 물리 장치, 수학 모델, 또는 계산 규칙을 뜻한다.
입력(input)은 밖에서 넣는 자극이고, 출력(output)은 우리가 관찰하기로 고른 결과이다.
상태(state)는 시스템 안에 남아 있는 기억이다.
예를 들어 질량-스프링-댐퍼에서는 외력이 입력이 될 수 있고, 위치가 출력이 될 수 있으며, 위치와 속도는 다음 움직임을 정하는 상태가 된다.
RLC 회로에서는 입력 전압을 넣고 커패시터 전압을 출력으로 볼 수 있으며, 전하와 전류가 회로의 기억 역할을 한다.

여기서 상태, 흐름, 출력은 서로 배타적인 물체 이름이 아니라 읽는 역할이다.
상태는 다음 순간을 계산하기 위해 시스템 안에 남아 있어야 하는 기억이고, 흐름은 전류 $i$나 속도 $\dot{x}$처럼 에너지와 전력이 지나가는 비율이며, 출력은 우리가 관찰하기로 고른 신호이다.
따라서 같은 전류도 회로의 상태 변수로 쓰일 수 있고, 동시에 전기적 흐름 변수로 읽힐 수 있다.

LTI 시스템은 입력과 출력의 관계가 선형이고, 시간이 지나도 그 관계가 변하지 않는 시스템이다.
이 장이 필요한 이유는 단순하다.
전기회로, 질량-스프링-댐퍼, 로봇 끝단 임피던스는 모두 시간에 따라 변하는 신호를 다룬다.
전압, 전류, 힘, 위치, 속도는 한 순간의 숫자가 아니라 시간 함수이다.
따라서 입력을 넣으면 출력이 나온다는 말을 정량적으로 다루려면, 숫자 하나가 아니라 시간 함수 전체를 다루는 언어가 필요하다.
LTI 시스템은 그중 가장 다루기 쉬우면서도 공학적으로 강력한 언어이다.

물론 실제 로봇은 완전히 선형도 아니고 완전히 시불변도 아니다.
마찰은 복잡하고, 관절 한계가 있으며, 큰 변위에서는 자코비안과 관성이 변한다.
그럼에도 LTI 모델을 먼저 배우는 이유는 작은 움직임, 한 작동점 근처, 교육용 데모에서는 많은 현상을 충분히 잘 설명하기 때문이다.
무엇보다 LTI 모델은 복잡한 시스템을 고유진동수, 감쇠비, 전달함수, 주파수 응답 같은 직관적인 도구로 나누어 볼 수 있게 해준다.

이 장에서는 네 가지 질문을 차례로 붙잡는다.
첫째, 어떤 관계를 선형이라고 부르는가.
둘째, 같은 실험을 나중에 해도 같은 결과가 나온다는 말은 수식으로 무엇인가.
셋째, 왜 라플라스 변환과 전달함수가 미분방정식을 읽기 쉽게 만드는가.
넷째, 왜 정적 정상상태에서는 스프링처럼 보이던 시스템도 빠르게 움직이면 질량과 감쇠의 영향을 강하게 받는가.
이 네 질문이 이어지면 뒤의 RLC 회로, 매스-스프링-댐퍼, 로봇 임피던스 제어가 같은 2차 시스템 언어로 연결된다.
앞 장의 통신선 이야기가 크기와 위상이 시간에 따라 어떻게 바뀌는가라는 질문을 남겼다면, 이 장은 그 질문을 입력, 출력, 상태, 전달함수라는 문법으로 바꾸는 단계이다.

그 연결을 미리 한 장으로 보면 다음과 같다.
직렬 RLC 회로를 전하 $q$로 쓰면
\begin{equation}
  L\ddot{q}+R\dot{q}+\frac{1}{C}q=v_{\mathrm{in}}
\end{equation}
이고, 매스-스프링-댐퍼를 위치 $x$로 쓰면
\begin{equation}
  m\ddot{x}+b\dot{x}+kx=f
\end{equation}
이다.
두 식은 모양이 거의 같다.
다만 전기에서는 전하와 전류가 움직이고, 기계에서는 위치와 속도가 움직인다.

\begin{table}[!htbp]
  \centering
  \caption{RLC 회로와 매스-스프링-댐퍼를 LTI 관점에서 나란히 보기}
  \small
  \setlength{\tabcolsep}{4pt}
  \begin{tabularx}{\linewidth}{@{}>{\raggedright\arraybackslash}p{0.19\linewidth}>{\raggedright\arraybackslash}p{0.22\linewidth}>{\raggedright\arraybackslash}p{0.22\linewidth}>{\raggedright\arraybackslash}X@{}}
    \toprule
    관점 & 전기 RLC & 기계 MSD & 직관 \\
    \midrule
    변위처럼 쌓이는 상태 & 전하 $q$ & 위치 $x$ & 얼마나 축적되었는가 \\
    흐름 상태 & 전류 $i=\dot{q}$ & 속도 $\dot{x}$ & 얼마나 빠르게 변하는가 \\
    변화 저항/저장 & 인덕턴스 $L$ & 질량 $m$ & 흐름의 급격한 변화를 싫어하고 에너지를 저장한다 \\
    소산 & 저항 $R$ & 감쇠계수 $b$ & 에너지를 열이나 마찰로 잃게 한다 \\
    복원 & 역커패시턴스 $1/C$ & 강성 $k$ & 기준 상태에서 벗어난 양을 되돌리려 한다 \\
    외부 자극 & 입력 전압 $v_{\mathrm{in}}$ & 외력 $f$ & 시스템을 움직이게 하는 원인 \\
    \bottomrule
  \end{tabularx}
\end{table}

이 표는 완벽한 일대일 번역표라기보다 이 글에서 사용할 교육용 지도에 가깝다.
전기 회로와 기계 시스템은 물리적으로 다르지만, 작은 신호와 LTI 가정 아래에서는 같은 수학 구조를 공유한다.

\subsection{선형성}

시스템을 연산자 $\mathcal{S}$로 나타내고, 입력 $u(t)$에 대한 출력을 $y(t)$라고 하자.
\begin{equation}
  y(t) = \mathcal{S}\{u(t)\}.
\end{equation}
여기서 $\mathcal{S}$는 숫자 하나를 숫자 하나로 바꾸는 함수라기보다, 입력 신호 전체를 출력 신호 전체로 바꾸는 규칙이다.
더 쉽게 말하면 입력 그래프 하나를 받아 출력 그래프 하나를 만들어 내는 규칙이다.
예를 들어 질량-스프링-댐퍼에서는 입력 $u(t)=f(t)$, 출력 $y(t)=x(t)$로 잡을 수 있다.
RLC 회로에서는 입력 $u(t)=v_{\mathrm{in}}(t)$, 출력 $y(t)=q(t)$ 또는 $v_C(t)$로 잡을 수 있다.
같은 물리 장치라도 무엇을 입력과 출력으로 고르느냐에 따라 서로 다른 전달함수로 보인다.
시스템이 선형이라는 것은 임의의 입력 $u_1(t)$, $u_2(t)$와 스칼라 $a$, $b$에 대해 다음이 성립한다는 뜻이다.
\begin{equation}
  \mathcal{S}\{a u_1(t) + b u_2(t)\}
  =
  a \mathcal{S}\{u_1(t)\}
  +
  b \mathcal{S}\{u_2(t)\}.
\end{equation}
이 식은 가산성(additivity)과 비례성(proportionality)을 하나로 합친 표현이다.
즉, 입력을 더한 뒤 시스템에 넣어도 되고, 각각의 출력을 구한 뒤 더해도 결과가 같다.
또한 입력을 두 배로 만들면 출력도 두 배가 된다.
보통 전달함수 문맥에서의 선형성은 초기조건을 0으로 둔 입력-출력 관계, 또는 평형점 주변의 작은 변화량에 대한 관계를 말한다.
여기서 말하는 선형성은 학교 수학에서 말하는 직선 그래프와 조금 다르다.
zero-state 입력-출력 관계에서 선형 시스템은 입력이 0이면 출력도 0이어야 한다.
초기 에너지가 남아 있으면 입력이 없어도 zero-input 응답이 생길 수 있으므로, 이 문장은 초기조건을 0으로 둔 관계에 대한 말이다.
따라서 $y=mx+b$처럼 원점을 지나지 않는 관계는 엄밀히는 affine, 즉 선형항에 상수 오프셋이 더해진 관계이다.
첫 단계에서는 이 차이를 단순하게 기억하면 된다.
직선처럼 보이더라도 원점을 지나지 않으면 제어에서 말하는 선형 시스템은 아니고, 선형식에 바탕값이 더해진 관계이다.
원점을 지나는 비례 관계는 선형이고, 일정한 바탕값이 더해져 있으면 그 자체로는 엄밀한 의미의 선형이 아니다.

두 번째 단계에서는 오프셋을 빼고 볼 수 있다.
평형점 $x_0$ 주변의 변화량 $\tilde{x}=x-x_0$를 쓰면, 그 위치에 스프링을 붙잡아 두기 위해 외부에서 가해야 하는 유지힘은 $f_{\mathrm{hold}}=k(x-x_0)=k\tilde{x}$처럼 선형식으로 볼 수 있다.
반대로 스프링이 물체에 실제로 가하는 복원력은 $f_s=-k\tilde{x}$이다.
로봇의 중력 보상이나 작동점 근처 선형화도 이런 식으로 오프셋을 떼어 내고 작은 변화량을 보는 과정이다.

예를 들어 어떤 스프링에 $1~\mathrm{N}$을 가했을 때 $1~\mathrm{cm}$ 늘어나고, $2~\mathrm{N}$을 가했을 때 $2~\mathrm{cm}$ 늘어난다면 그 범위에서는 선형이라고 볼 수 있다.
$1~\mathrm{N}$과 $2~\mathrm{N}$을 따로 생각한 뒤 변형을 더하면 $3~\mathrm{cm}$이고, 실제로 $3~\mathrm{N}$을 한 번에 가했을 때도 $3~\mathrm{cm}$가 나오기 때문이다.
또 $1~\mathrm{N}$에 대한 응답을 세 배 하면 $3~\mathrm{N}$에 대한 응답과 같다.
반대로 고무줄을 너무 많이 늘려서 갑자기 딱딱해지거나, 재료가 영구 변형되거나, 마찰 때문에 움직이기 전까지는 아무 변화가 없다면 선형성이 깨진다.
선형성은 자연의 모든 현상이 원래 선형이라는 뜻이 아니라, 특정 범위에서 입력과 출력의 관계를 직선처럼 볼 수 있다는 뜻이다.

\subsection{시불변성}

시불변성은 시스템의 성질이 시간에 따라 변하지 않는다는 뜻이다.
입력 $u(t)$에 대한 출력이 $y(t)$라면, 같은 입력을 $t_0$만큼 늦게 넣었을 때 출력도 똑같이 $t_0$만큼 늦어져야 한다.
\begin{equation}
  \mathcal{S}\{u(t-t_0)\} = y(t-t_0).
\end{equation}
예를 들어 같은 질량-스프링-댐퍼 시스템을 오전에 실험하든 오후에 실험하든, 물리 파라미터가 변하지 않는다면 같은 입력에 대해 같은 형태의 응답을 보여야 한다.
조금 더 그래프처럼 말하면, 0초에 짧은 힘 펄스를 넣었을 때 0.3초 뒤 최대 변위가 나온다면, 5초에 같은 펄스를 넣었을 때는 5.3초 근처에서 같은 모양의 최대 변위가 나와야 한다.
단, zero-state 응답, 즉 입력이 들어오기 전 저장 에너지와 초기 운동이 없는 응답을 비교하려면 실험을 시작할 때 시스템은 같은 초기상태로 리셋되어 있어야 한다.
반대로 시간이 지나며 감쇠기가 뜨거워져 $b(t)$가 변하거나, 제어기가 중간에 $K_p(t)$를 바꾸거나, 배터리 전압이 떨어져 같은 명령에도 힘이 달라지면 시불변성이 깨진다.

\subsection{전달함수}

전달함수로 넘어가기 전에는 세 가지를 먼저 확인하자.
입력과 출력이 무엇인지 정했는가?
초기 에너지를 0으로 두고 보는가?
지금 보는 범위에서 선형성과 시불변성을 가정할 수 있는가?
이 조건이 정해져야 $H(s)=Y(s)/U(s)$가 우연히 두 파형을 나눈 값이 아니라, 시스템을 읽는 약속된 방법이 된다.

그다음에는 응답의 종류를 구분해 두자.
처음에는 이 단어들이 비슷하게 들리지만, 무엇이 시스템을 움직이게 했는지를 나누는 이름일 뿐이다.

\begin{table}[!htbp]
  \centering
  \caption{전달함수를 읽기 전에 구분할 응답의 종류}
  \small
  \setlength{\tabcolsep}{4pt}
  \begin{tabularx}{\linewidth}{@{}>{\raggedright\arraybackslash}p{0.22\linewidth}>{\raggedright\arraybackslash}p{0.29\linewidth}>{\raggedright\arraybackslash}X@{}}
    \toprule
    이름 & 무엇이 움직임을 만들었는가 & 예시 \\
    \midrule
    zero-input 응답 & 입력은 없지만 초기 위치, 초기 속도, 저장 전하처럼 처음부터 에너지가 있음 & 당겨 놓은 스프링을 그냥 놓았을 때의 자유진동 \\
    zero-state 응답 & 초기 에너지는 없고, 새로 들어온 입력만 응답을 만듦 & 정지해 있던 물체를 외력으로 밀기 시작한 경우 \\
    전체 응답 & zero-input 응답과 zero-state 응답이 함께 나타남 & 이미 흔들리던 물체에 외력을 추가로 가한 경우 \\
    임펄스 응답 & 아주 짧고 강한 한 번의 두드림에 대한 zero-state 응답 & 책상 위 물체를 순간적으로 톡 친 뒤의 움직임 \\
    계단 응답 & 일정한 입력을 어느 순간부터 계속 가했을 때의 응답 & $t=0$부터 일정한 힘 $F_0$로 계속 미는 경우 \\
    \bottomrule
  \end{tabularx}
\end{table}

전달함수는 이 중에서 zero-state 응답을 다룬다.
초기조건을 0으로 둔다는 것은 입력을 넣기 전에는 이미 움직이던 속도나 저장된 에너지가 없다고 보는 것이다.
이미 흔들리던 질량이나 충전된 커패시터는 입력이 없어도 자기 응답을 만든다.
그 응답까지 입력-출력 비율에 섞이면, 시스템 자체가 새 입력을 어떻게 바꾸는지 흐려진다.

전달함수를 만들 때는 매번 세 가지를 먼저 고른다.
무엇을 입력으로 넣는가, 무엇을 출력으로 볼 것인가, 그리고 처음에 저장된 에너지를 0으로 둘 것인가이다.
라플라스 변환의 역할도 여기서 나온다.
시간영역의 미분방정식은 말 그대로 시간에 따라 계속 변하는 값을 따라가야 한다.
예를 들어
\begin{equation}
  m\ddot{x}+b\dot{x}+kx=f
\end{equation}
를 시간영역에서 직접 풀면 위치, 속도, 가속도의 관계를 계속 추적해야 한다.
라플라스 변환은 이런 미분 관계를 $s$라는 복소 변수의 대수식으로 바꾸어 준다.
시간을 없애는 마법이라기보다, 시간에 따라 계속 갱신해야 하는 미분 장부를 한 장의 대수 장부로 바꾸는 도구에 가깝다.
이렇게 바꾸면 같은 시스템에 다른 입력을 넣어도 매번 운동방정식을 처음부터 풀지 않고, $H(s)$라는 입력-출력 변환표를 재사용할 수 있다.
$X(s)$는 $x(t)$를 라플라스 변환한 것이고, $F(s)$는 $f(t)$를 라플라스 변환한 것이다.

\paragraph{라플라스 변환의 정의.}
이 변환이 정확히 무엇인지 한 번은 정의로 적어 두자.
\begin{equation}
  X(s)
  =
  \mathcal{L}\{x(t)\}
  =
  \int_{0}^{\infty} x(t)\,e^{-st}\,\dd t.
  \label{eq:laplace-definition}
\vmark{laplace-definition-integral}%
\end{equation}
말로 풀면, 시간 함수 $x(t)$에 $e^{-st}$라는 점점 작아지는 추를 곱해서 $t=0$부터 끝까지 모두 더한(적분한) 값이다.
적분 기호 $\int$는 아주 잘게 나눈 조각 값들의 합으로 읽으면 된다.
$s$를 하나 고를 때마다 숫자가 하나 나오므로, 시간 파형 전체가 $s$의 함수 $X(s)$ 한 장으로 요약된다.

\paragraph{미분 규칙은 곱의 미분에서 나온다.}
아주 거칠게 말하면, 미분은 $s$를 곱하는 연산처럼 바뀐다.
이 규칙은 외워야 하는 별개의 공식이 아니라, 위 정의와 곱의 미분 규칙만으로 나온다.
한 단계씩 보자.
$s$는 $t$와 무관한 상수 취급이므로, 곱의 미분 규칙에 따라
\[
  \frac{\dd}{\dd t}\left[x(t)\,e^{-st}\right]
  =
  \dot{x}(t)\,e^{-st}
  -
  s\,x(t)\,e^{-st}
\vmark{laplace-rule-product-rule-derivation}%
\]
이다.
양변을 $t=0$부터 $\infty$까지 적분하자.
왼쪽은 미분한 것을 다시 적분하는 것이므로 끝값 빼기 시작값, 즉 $x(t)e^{-st}$의 $t\to\infty$ 값에서 $t=0$ 값 $x(0)$을 뺀 것이 된다.
응답이 발산하지 않아 $t\to\infty$에서 $x(t)e^{-st}\to0$이 되는 $s$에서 보면 왼쪽은 $0-x(0)=-x(0)$이다.
오른쪽의 두 적분은 정의 \eqref{eq:laplace-definition}에 따라 각각 $\mathcal{L}\{\dot{x}\}$와 $-sX(s)$이다.
따라서 $-x(0)=\mathcal{L}\{\dot{x}\}-sX(s)$이고, 정리하면
\begin{equation}
  \mathcal{L}\{\dot{x}(t)\}=sX(s)-x(0)
\end{equation}
이고, 초기조건을 0으로 두면 $\dot{x}$는 $sX(s)$처럼 다룰 수 있다.
간단한 검산도 해 보자.
$x(t)=1$(상수)이면 정의에서 $X(s)=\int_0^\infty e^{-st}\dd t=\left[-e^{-st}/s\right]_0^\infty=0-\left(-\tfrac{1}{s}\right)=\tfrac{1}{s}$이다.
$\dot{x}=0$이므로 규칙의 왼쪽은 $\mathcal{L}\{0\}=0$이고, 오른쪽도 $sX(s)-x(0)=s\cdot(1/s)-1=0$으로 일치한다.
\vmark{laplace-rule-constant-check}%

두 번 미분한 항은 같은 규칙을 두 번 적용하면 된다.
$\dot{x}$를 새 함수로 보고 규칙을 다시 쓰면
\[
  \mathcal{L}\{\ddot{x}\}
  =
  s\,\mathcal{L}\{\dot{x}\}-\dot{x}(0)
  =
  s^2X(s)-s\,x(0)-\dot{x}(0)
\vmark{laplace-second-derivative-twice}%
\]
이고, 초기조건을 0으로 두면 $\ddot{x}$는 $s^2X(s)$처럼 다룰 수 있다.
이제 $m\ddot{x}+b\dot{x}+kx=f$의 각 항을 하나씩 바꾸면 $m\ddot{x}\to ms^2X(s)$, $b\dot{x}\to bsX(s)$, $kx\to kX(s)$, $f\to F(s)$이므로
\vmark{laplace-term-by-term-msd}%
\begin{equation}
  (ms^2+bs+k)X(s)=F(s)
\end{equation}
가 된다.
그래서 미분방정식이 다항식 비율로 정리된다.

여기서 $s$는 단순한 알파벳이 아니라, 시간응답의 성질을 담는 복소 변수이다.
실제로 회로에 넣거나 로봇에 명령하는 물리 입력이 아니라, 이런 모양의 응답이 가능하다면 어떻게 보일까를 묻는 시험용 주소라고 생각하면 된다.
보통
\begin{equation}
  s=\sigma+\jj\omega
\end{equation}
로 생각한다.
$\sigma$는 응답이 커지거나 줄어드는 속도와 관련되고, $\omega$는 흔들리는 속도와 관련된다.
실제 물리 입력이 복소수라는 뜻은 아니다.
감쇠와 진동을 한 기호 안에서 함께 계산하기 위한 표현이다.
특히 주파수 응답을 볼 때는 $s=\jj\omega$를 대입한다.
안정하고 $\jj\omega$가 시스템의 pole, 즉 전달함수 분모를 0으로 만드는 특이한 값이 아닐 때, 이 대입은 실제 정현파 정상상태의 크기와 위상을 읽는 방법이 된다.
pole은 단순히 수학적으로 분모가 0이 되는 점이 아니라, 시스템이 자연스럽게 가지고 있는 감쇠 모드와 진동 모드를 정하는 값이다.
이 주소들 중 어떤 곳에서는 외부에서 계속 밀지 않아도 시스템 고유의 자유응답 모양이 드러난다.
그래서 pole의 위치를 보면 응답이 발산할지, 진동할지, 얼마나 빨리 사라질지 알 수 있다.
고유진동수, 감쇠비, 오버슈트, 정착시간은 모두 이런 응답 모양을 읽기 위한 이름이다.
이 장에서는 pole, 컨볼루션, 최종값 정리를 모두 완벽히 다루려 하지 않는다.
지금 꼭 붙잡을 것은 입력-출력 비율과 응답 모양이 전달함수의 분모 구조에서 나온다는 점이고, 각 용어의 물리적 의미는 뒤의 기계 시스템 장에서 다시 풀어 본다.
제어공학에서 라플라스 변환이 중요한 이유는 계산이 편하기 때문만이 아니라, 응답 모양을 읽는 언어가 되기 때문이다.

이 절에서는 입력이 들어오기 전에는 시스템이 정지해 있고, 원인보다 결과가 먼저 나오지 않는 causal LTI 시스템을 생각한다.
causal은 원인보다 결과가 먼저 나타나지 않는다는 뜻이다.
즉, 초기조건이 0이고 $t<0$에서는 $u(t)=0$인 zero-state 응답을 기본으로 본다.
LTI 시스템은 임펄스 응답 $h(t)$와 컨볼루션으로 표현할 수 있다.
임펄스 응답은 시스템을 아주 짧고 강하게 한 번 두드렸을 때 나오는 반응이다.
컨볼루션은 임의의 입력을 아주 작은 두드림들의 합으로 나누고, 각각의 반응을 시간에 맞춰 더하는 계산이라고 이해하면 된다.
예를 들어 0.2초에 작은 힘을 한 번 주면 그 뒤의 반응은 $h(t-0.2)$ 모양으로 나타난다.
0.5초에 두 배 세게 밀면 $2h(t-0.5)$가 더해진다.
책상 위 물체를 여러 번 톡톡 치면, 각 두드림에 대한 작은 흔들림들이 시간축 위에서 겹쳐 하나의 전체 움직임을 만든다.
실제 입력은 이런 아주 작은 밀기들이 촘촘히 이어진 것으로 보고, 각각의 반응을 시간에 맞춰 모두 더한 것이 컨볼루션이다.
처음 읽을 때 이 적분을 완전히 계산할 필요는 없다.
중요한 것은 LTI 시스템에서는 한 번 두드렸을 때의 반응만 알면, 여러 입력에 대한 반응을 조립할 수 있다는 점이다.
이 글처럼 causal이고, $t<0$에서 입력과 임펄스 응답이 0인 zero-state 시스템에서는 $t\ge0$에 대해
\begin{equation}
  y(t) = (h * u)(t) = \int_{0}^{t} h(t-\tau) u(\tau) \, \dd \tau
\end{equation}
로 쓸 수 있다.
일반적인 컨볼루션 표기에서는 적분 범위를 $-\infty$부터 $\infty$까지 쓰기도 하지만, 여기서는 물리적으로 입력을 시작하는 시각을 $t=0$으로 잡은 인과 시스템을 다루기 때문에 위처럼 적었다.

전달함수(transfer function)는 zero-state 임펄스 응답 $h(t)$의 라플라스 변환이다.
\begin{equation}
  H(s)=\mathcal{L}\{h(t)\}.
\end{equation}
라플라스 영역에서는 컨볼루션이 곱셈으로 바뀌기 때문에
\begin{equation}
  Y(s) = H(s) U(s)
\end{equation}
가 된다.
따라서 같은 조건에서 입력과 출력의 비율로 쓰면
\begin{equation}
  H(s) = \frac{Y(s)}{U(s)}
\end{equation}
이다.
정확히 말하면 전달함수는 SISO, 즉 입력 하나와 출력 하나를 보는 LTI 시스템의 zero-state 입력-출력 연산자이다.
그래서 전달함수는 특정 입력 하나를 우연히 나눈 값이 아니라, 그 시스템이 입력 신호를 출력 신호로 바꾸는 방식 자체를 담는다.

전달함수를 만들 때의 순서도 매번 같다.
먼저 에너지가 쌓이는 대표 좌표를 고르고, 다음에 외부에서 넣는 입력을 정한다.
그다음 우리가 보고 싶은 출력을 고른 뒤, 초기 저장 에너지는 0이라고 두고 $Y(s)/U(s)$를 계산한다.
이 순서를 지키면 같은 물리 장치에서도 왜 위치 출력, 속도 출력, 전류 출력의 전달함수가 달라지는지 구분이 선명해진다.

무엇을 입력과 출력으로 잡느냐에 따라 같은 물리 시스템도 다른 전달함수로 보일 수 있다.
기계 시스템에서는 특히 다음 용어를 구분해야 한다.
아래 표에서 $V_x(s)$는 속도 $\dot{x}(t)$의 라플라스 변환이다.
zero-state 조건에서는 $V_x(s)=sX(s)$로 볼 수 있다.

\begin{table}[!htbp]
  \centering
  \caption{기계 시스템에서 입력과 출력을 어떻게 고르느냐에 따른 이름과 단위}
  \small
  \setlength{\tabcolsep}{4pt}
  \begin{tabularx}{\linewidth}{@{}>{\raggedright\arraybackslash}p{0.22\linewidth}>{\raggedright\arraybackslash}p{0.20\linewidth}>{\raggedright\arraybackslash}p{0.18\linewidth}>{\raggedright\arraybackslash}X@{}}
    \toprule
    이름 & 비율 & 단위 & 의미 \\
    \midrule
    순응성(compliance) & $X(s)/F(s)$ & m/N & 힘을 받았을 때 얼마나 움직이는가 \\
    어드미턴스 또는 모빌리티 & $V_x(s)/F(s)$ & (m/s)/N & 힘을 받았을 때 얼마나 빠르게 움직이는가 \\
    기계 임피던스 & $F(s)/V_x(s)$ & N\,s/m & 어떤 속도 운동을 만들기 위해 얼마나 큰 힘이 필요한가 \\
    동적 강성(dynamic stiffness) & $F(s)/X(s)$ & N/m & 어떤 변위를 만들기 위해 얼마나 큰 힘이 필요한가 \\
    \bottomrule
  \end{tabularx}
\end{table}

전기 임피던스는 보통 $Z(s)=V(s)/I(s)$처럼 effort를 flow로 나눈다.
기계 임피던스도 엄밀한 주파수영역 정의에서는 $Z_m(s)=F(s)/V_x(s)$처럼 힘을 속도로 나눈다.
다만 로봇 임피던스 제어에서는 원하는 질량-스프링-댐퍼 관계 자체를 설계한다는 뜻으로 임피던스라는 말을 넓게 쓴다.
그래서 정지한 접촉 상황에서는 $F/X$인 강성 또는 그 역수인 순응성을 주로 보게 되고, 빠른 운동이나 진동을 보면 $F/V_x$인 기계 임피던스 관점이 중요해진다.

예를 들어 질량-스프링-댐퍼에서 입력을 힘, 출력을 위치로 잡으면
\begin{equation}
  H(s)=\frac{X(s)}{F(s)}
  =
  \frac{1}{ms^2+bs+k}
\end{equation}
이다.
분모의 $ms^2$ 항은 가속도와 관성의 영향을, $bs$ 항은 속도와 감쇠의 영향을, $k$ 항은 위치와 스프링 복원의 영향을 담는다.
$H(0)=1/k$는 정적 순응성(static compliance)이며 단위는 $\mathrm{m/N}$이다.
따라서 안정하고 $k>0$인 시스템에 일정한 힘 $F_0$를 오래 가하면 최종 변위는
\begin{equation}
  x_\infty = \frac{F_0}{k}
\end{equation}
가 된다.
여기서 $H(0)$ 자체가 최종 변위인 것은 아니다.
단위 힘에 대해 얼마나 움직이는지를 나타내는 비율이고, 실제 최종 변위는 입력 힘의 크기를 곱해야 한다.
분모 $ms^2+bs+k$의 근, 즉 pole은 질량, 감쇠, 스프링이 함께 만드는 자연스러운 응답 모양을 담고, 그 위치가 진동 여부와 감쇠 속도를 정한다.

\subsection{정현파 입력과 주파수 응답}

LTI 시스템의 중요한 성질은 정현파 입력에 대한 응답이다.
입력이
\begin{equation}
  u(t) = A\cos(\omega t)
\end{equation}
처럼 하나의 주파수 $\omega$를 갖는 정현파라면, 안정하고 $\jj\omega$에 pole이 없는 LTI 시스템의 정상상태 출력은 같은 주파수를 가진다.
달라지는 것은 크기와 위상뿐이다.
\begin{equation}
  y_{\mathrm{ss}}(t)
  =
  A|H(\jj\omega)|
  \cos\left(\omega t + \angle H(\jj\omega)\right).
\end{equation}
여기서 $|H(\jj\omega)|$와 $\angle H(\jj\omega)$는 복소수 $H(\jj\omega)$의 크기와 각도이다.
복소수 $a+\jj b$를 평면 위의 점 $(a,b)$로 그리면, 크기 $|a+\jj b|=\sqrt{a^2+b^2}$는 원점에서 그 점까지의 거리이고, 각도 $\angle(a+\jj b)$는 양의 실수축에서 그 점 방향까지 잰 회전각이다.
예를 들어 어떤 주파수에서 $H(\jj\omega)=1-\jj$라면 크기는 $\sqrt{1^2+(-1)^2}=\sqrt{2}$, 각도는 $-45^\circ$이므로, 출력 진폭은 입력의 약 1.41배가 되고 위상은 $45^\circ$만큼 늦는다.
\vmark{freq-magnitude-angle-complex}%
따라서 $H(\jj\omega)$는 입력 주파수에 대해 시스템이 얼마나 증폭하거나 감쇠하는지, 그리고 얼마나 늦거나 빠르게 반응하는지를 알려준다.
예를 들어 어떤 주파수에서 $|H(\jj\omega)|=0.5$라면 출력 진폭은 입력 진폭의 절반이다.
$\angle H(\jj\omega)=-90^\circ$라면 출력은 같은 주파수로 흔들리지만 입력보다 한 주기의 1/4만큼 늦게 나타난다.
이처럼 주파수 응답은 같은 입력이라도 천천히 밀 때와 빠르게 흔들 때 시스템이 다르게 느껴지는 이유를 설명한다.
천천히 밀 때는 $s\approx0$에 가까워 스프링 같은 정적 관계가 두드러지고, 빠르게 흔들 때는 질량과 감쇠 때문에 크기와 위상 차이가 크게 달라진다.

임피던스도 같은 방식으로 해석된다.
전기 임피던스 $Z(\jj\omega)$는 특정 주파수에서 전류에 비해 전압이 얼마나 필요한지, 그리고 전압과 전류 사이의 위상차가 얼마인지를 나타낸다.
기계 임피던스 $Z_m(\jj\omega)$는 특정 주파수의 속도 운동을 만들기 위해 얼마나 큰 힘이 필요한지, 그리고 힘과 속도 사이에 어떤 위상차가 생기는지를 나타낸다.
이 때문에 임피던스는 단순한 상수가 아니라 주파수에 따라 달라지는 동역학적 저항으로 이해해야 한다.

\subsection{정적 정상상태응답과 과도응답}

제어와 임피던스를 읽을 때는 정상상태응답과 과도응답을 구분해야 한다.
여기서 먼저 말하는 것은 일정한 힘이나 직류 입력을 오래 가했을 때의 정적 정상상태응답(static steady-state response)이다.
앞 절의 정현파 정상상태에서는 질량과 감쇠도 크기와 위상에 계속 영향을 주지만, 정적 정상상태에서는 최종적으로 움직임이 멈춘다.
예를 들어 일정한 힘으로 스프링을 천천히 밀고 오래 기다리면, 물체는 더 이상 움직이지 않는 위치에 머문다.
이때 속도와 가속도는 0이고, 남는 것은 스프링 힘과 외력의 균형이다.

과도응답(transient response)은 그 최종 상태에 도달하기 전까지의 과정이다.
힘을 갑자기 가하면 물체는 바로 최종 위치로 순간이동하지 않는다.
질량 때문에 가속도가 생기고, 속도가 생기고, 감쇠가 부족하면 목표점을 지나쳐 흔들린다.
따라서 정적 정상상태에서는 스프링 항만 남는 경우가 많지만, 과도응답에서는 질량과 감쇠가 반드시 중요해진다.

임피던스 제어에서 이 구분은 매우 실용적이다.
접촉 중 최종 힘과 최종 변위를 설명할 때는 크기 관계로 $|f|\approx k|x|$라는 스프링 관계만으로도 충분할 수 있다.
부호까지 쓰려면 먼저 힘과 좌표의 기준을 정해야 한다.
여기서 $x$를 외력이 미는 방향의 변위로 잡고, $f_{\mathrm{ext}}$를 질량이나 로봇에 가해진 외력으로 정의하면 정적 평형에서 $f_{\mathrm{ext}}\approx kx$이고, 스프링 또는 제어기가 물체에 가하는 복원력은 $f_s\approx-kx$가 된다.
반대로 벽 침투량 $\delta=x-x_{\mathrm{wall}}>0$을 벽 안쪽 양의 방향으로 잡으면, 벽이 로봇에 가하는 힘은 보통 $f_{\mathrm{wall}}\approx-k\delta$처럼 반대 부호로 쓴다.
하지만 로봇이 벽에 닿는 순간의 충격, 목표점으로 도달하는 속도, 오버슈트, 진동, 정착시간을 설명하려면 $m$, $b$, $k$가 모두 들어간 2차 동역학을 봐야 한다.
이 글이 LTI 시스템을 거쳐 2차 전기 시스템과 2차 기계 시스템을 다루는 이유가 바로 여기에 있다.

다음 장의 자유진동은 입력은 0이지만 초기 에너지가 있는 경우이고, 전달함수 계산은 초기 에너지가 0이고 입력이 있는 경우이다.
둘 다 같은 LTI 방정식에서 나오지만, 무엇이 움직임을 시작시켰는지가 다르다.

나중에 사용할 최종값 정리(final value theorem)도 같은 생각에서 나온다.
이 정리는 응답이 최종적으로 한 값에 안정되게 수렴할 때만 쓸 수 있다.
계속 진동하거나 발산하거나 일정한 속도로 계속 커지는 응답에는 쓰면 안 된다.
시간이 충분히 지난 뒤의 최종값을 알고 싶을 때, 최종값이 실제로 존재하고 $sY(s)$의 모든 pole이 열린 좌반평면에 있으면 시간응답을 끝까지 풀지 않고도
\begin{equation}
  \lim_{t\to\infty} y(t)
  =
  \lim_{s\to 0} sY(s)
\end{equation}
로 계산할 수 있다.
좌반평면은 복소평면에서 실수부가 음수인 영역이고, 물리적으로는 시간이 지나며 사라지는 모드에 해당한다.
우반평면 pole, 허축 pole, 감춰진 불안정 pole-zero cancellation, $sY(s)$에 원점 pole이 남는 적분 동작, 감쇠 없는 지속 진동이 있으면 이 정리를 조심해야 한다.
계단 입력 때문에 $Y(s)$에 생긴 단순 원점 pole은 $sY(s)$에서 제거된 뒤를 본다고 이해하면 된다.
처음에는 pole 조건을 모두 판별하려고 멈춰 서기보다, 최종값 정리는 응답이 실제로 한 값에 멈출 때만 쓰는 도구라고 잡아 두자.
계단 입력의 $1/s$는 일정한 입력을 계속 붙잡고 있다는 표시이고, $sY(s)$에서 $s\to0$을 보는 것은 빠른 흔들림이 모두 사라진 뒤 남는 직류 균형만 묻는 것과 비슷하다.
예를 들어 감쇠가 있는 질량-스프링-댐퍼에 일정한 힘을 주면 결국 한 위치에서 멈추므로 최종값 정리를 쓸 수 있다.
반대로 감쇠가 없는 스프링-질량계를 당겨 놓고 놓으면 계속 흔들리므로 최종 위치 하나가 존재하지 않는다.
모터가 일정한 속도로 계속 움직이는 적분기형 위치 출력도 시간이 지날수록 위치가 계속 커지므로 최종 위치를 하나로 정할 수 없다.
직관적으로 $s\to 0$은 아주 느린 변화, 즉 직류나 최종 균형 상태를 보는 것과 비슷하다.
그래서 일정한 외력에 대한 최종 위치 오차를 구할 때 질량과 감쇠항은 사라지고 스프링 강성만 남는다.

예를 들어 힘 입력이 $f(t)=F_0$인 계단 입력이면 $F(s)=F_0/s$이고,
\begin{equation}
  X(s)=\frac{1}{ms^2+bs+k}\frac{F_0}{s}
\end{equation}
이다.
$m>0$, $b>0$, $k>0$로 점근 안정이고 최종값이 존재하면
\begin{equation}
  x_\infty
  =
  \lim_{s\to0}sX(s)
  =
  \frac{F_0}{k}
\end{equation}
가 된다.
이 계산은 정상상태에서는 속도와 가속도가 0이 되어 결국 스프링 힘과 외력이 균형을 이룬다는 물리 직관과 같다.

\subsection{LTI 관점에서 기억할 것}

이 장의 목적은 라플라스 변환을 능숙하게 계산하는 것이 아니라, 뒤의 전기 시스템과 기계 시스템을 같은 눈으로 볼 준비를 하는 것이다.
처음 읽을 때는 다음 정도만 확인해도 충분하다.

\begin{table}[!htbp]
  \centering
  \caption{뒤 장으로 넘어가기 전 LTI 체크리스트}
  \small
  \setlength{\tabcolsep}{4pt}
  \begin{tabularx}{\linewidth}{@{}>{\raggedright\arraybackslash}p{0.24\linewidth}>{\raggedright\arraybackslash}X@{}}
    \toprule
    질문 & 확인할 내용 \\
    \midrule
    무엇을 입력으로 잡았는가 & 힘, 전압, 전류, 목표 위치처럼 밖에서 넣는 신호가 무엇인지 정한다. \\
    무엇을 출력으로 볼 것인가 & 위치, 속도, 전류, 전압처럼 관찰하려는 신호를 정한다. 같은 장치도 출력 선택에 따라 전달함수가 달라진다. \\
    초기 에너지가 있는가 & 전달함수는 보통 zero-state 응답을 본다. 이미 흔들리거나 충전된 상태라면 zero-input 응답이 따로 생긴다. \\
    선형 근사가 가능한가 & 작은 움직임이나 평형점 근처에서는 선형으로 볼 수 있지만, 포화, 마찰, 큰 변위에서는 깨질 수 있다. \\
    시간이 지나도 파라미터가 같은가 & 질량, 감쇠, 강성, 저항, 제어 이득이 시간에 따라 바뀌면 시불변성이 깨진다. \\
    정적 최종값을 보는가, 과도응답을 보는가 & 정지한 최종 균형에서는 강성이 주로 보이고, 도달 속도와 진동에는 질량과 감쇠가 함께 보인다. \\
    \bottomrule
  \end{tabularx}
\end{table}

다음 장에서는 이 체크리스트를 먼저 RLC 회로에 적용하고, 그다음 질량-스프링-댐퍼와 로봇 끝단으로 옮겨 간다.
RLC에서는 $q$와 $i$를, 기계 시스템에서는 $x$와 $\dot{x}$를, 로봇에서는 위치 오차와 속도 오차를 같은 눈으로 읽게 된다.
인덕터, 저항, 커패시터가 각각 어떤 방식으로 전류와 전하의 변화를 막고 에너지를 저장하거나 소산하는지 보면, 임피던스가 단순한 저항보다 넓은 개념이라는 점이 더 분명해진다.
다음 장에서 풀 물리 질문은 이것이다.
저항, 인덕터, 커패시터가 같은 2차 방정식 안에서 각각 어떤 기억, 흐름, 소산 역할을 맡는가?
